% Generated by analysis/scripts/export_ten_cases.py
\subsection*{统一结果卡}
\begingroup\small
\noindent\begin{tabularx}{\textwidth}{@{}p{24mm}X@{}}\toprule
技术/模式 & PCM；SLC二状态8行电压式ACIM；32-IDAC对角写入\\
逻辑粒度/聚合 & B\_S=128 Byte；B\_R=32 Byte。32 Byte为对角选通，非任意连续32 Byte；d=0--127、s=0--3，512事务无重复覆盖矩阵\\
资源/相对R0 & 128 SAR、16重构通道、32实际IDAC与行/对角选通；8位平面、每向量1024轮\\
服务口径 & 无另列周期维护；完整波形、驱动预留与终验在resident内\\
主导因素/选择 & 参考SET/RESET占写时间73\%；8行组织增加读轮次使$\rho$及ridge较低，不能推出普适动态工作负载优势\\
关键对照 & 主二态终验与768 ns弱信号前端分开；SET总250/300 ns解释对照\\
\bottomrule\end{tabularx}
\vspace{3mm}
\begin{center}\begin{tabular}{@{}lrrr@{}}\toprule
成对条件情景 & $\rho$ (MB/s) & $\tau$ (MB/s) & $\mathrm{RI}^{*}$\\\midrule
短预算 & 5.207 & 7.972 & 0.6532\\
参考（推荐） & 2.5 & 6.874 & 0.3636\\
长预算 & 1.041 & 4.946 & 0.2106\\
\bottomrule\end{tabular}\end{center}
\noindent 参考原始占用：streaming 5.121e+04 ns；resident 4655 ns。
\par\noindent 范围为有限成对条件情景极值，不是物理不确定性包络。1 MB=$10^6$ Byte；整矩阵16384 Byte=16 KiB。源定位见本例证据笔记；公共基线shared\_baseline。
\endgroup
\clearpage
